\documentclass[a4paper, 11pt,oneside]{article}
\usepackage[
  top=1.5cm,
  bottom=1cm,
  left=2cm,
  right=1.5cm,
  headheight=25.22153pt, % as per the warning by fancyhdr
  includehead,includefoot,
  heightrounded, % to avoid spurious underfull messages
]{geometry}

\usepackage[T1]{fontenc}
\usepackage{microtype}
\usepackage{fancyhdr}
\usepackage{fancyvrb}
\usepackage{lipsum}
\usepackage{url}
\usepackage{listings}
\usepackage{lastpage}
\usepackage{enumitem}
\usepackage{datetime}
\usepackage{minted}

\settimeformat{hhmmsstime}
\yyyymmdddate

% Semester shown in the page header. Change it here.
\newcommand{\semester}{First Semester 2026-2027}

% Box used for the questions (same look as the other labs, but sized so
% that it does not stick out into the margin).
\newcommand{\qbox}[1]{%
  \noindent\fbox{\parbox{\dimexpr\linewidth-2\fboxsep-2\fboxrule\relax}{#1}}%
}

\pagestyle{fancy}
\fancyhf{} % clear all fields

\lhead{CMSC 125: Operating Systems \\ \semester}
\rhead{Institute of Computer Science \\ University of the Philippines Los Banos}
\rfoot{JACHermocilla | CC BY-NC-SA 4.0}
\cfoot{\thepage\ of \pageref{LastPage}}
\lfoot{Revision: \today\ \currenttime}
\renewcommand{\headrulewidth}{0.4pt}
\renewcommand{\footrulewidth}{0.4pt}

\begin{document}

\begin{center}
   {\LARGE \textbf{ICS-OS Lab 04: Lottery Scheduling}}
\end{center}

\section*{Objectives}
   At the end of this activity, you should be able to:
   \begin{enumerate}[itemsep=0pt,parsep=0pt]
       \item trace how ICS-OS selects the next process to run, from the timer interrupt down to the scheduler;
       \item explain the lottery scheduling algorithm (Waldspurger and Weihl, 1994) and how tickets achieve \textit{proportional-share} CPU allocation;
       \item implement lottery scheduling by modifying the ICS-OS default scheduler; and
       \item verify experimentally that the CPU time received by threads is proportional to their ticket holdings.
   \end{enumerate}

\section{Introduction}
The default ICS-OS scheduler is \textit{round-robin}: every runnable process gets an equal share of the CPU. In this lab you will replace it with a \textit{lottery scheduler}, which gives each process a number of \textbf{tickets} and holds a raffle every time the CPU must be handed to someone. A process holding three times the tickets of another wins about three times as often, and so receives about three times the CPU time.

\section{Prerequisites}

To proceed with this lab, you should have completed \texttt{Lab 01} and \texttt{Lab 02}. Build the kernel in the Docker build environment described in Lab 01 and the README. Most of the commands below are run relative to the \texttt{\$ICSOS\_HOME} directory. Update your local copy of the source code and create a new branch for this lab with the commands below.

\begin{minted}[frame=single,framesep=10pt]{bash}
$cd $ICSOS_HOME
$git checkout master
$git pull
$git checkout -b starter/lottery-scheduling-lab
$git branch    #to check the current branch
\end{minted}

Unless stated otherwise, all source paths in this document are under \texttt{ics-os/kernel/}.

\section{Deliverables and Credit}
Perform the tasks below and capture screen shots while you do them. Answer all questions. Submit a PDF file containing the screen shots with captions, your answers to the questions, and your modified \texttt{scheduler.c} (or the output of \texttt{git diff}). Do not forget to put your name and laboratory section. Total is \textbf{85 points}.

\section{Background}

\subsection*{How scheduling works in ICS-OS}
Study the following parts of the kernel (These have been discussed in the handout).

\begin{itemize}[itemsep=0pt,parsep=0pt]
  \item \texttt{process/process.h} defines \texttt{PCB386}, the process control block. Each PCB sits in a \textbf{circular doubly-linked list} (the \texttt{next} and \texttt{before} pointers). Note these fields: \texttt{status} (the \texttt{PS\_ATTB\_BLOCKED} bit marks a blocked process), \texttt{waiting} (a countdown for sleeping processes), \texttt{totalcputime} (incremented on every tick the process runs), and \texttt{priority}, which the default scheduler \textbf{never uses}. You will reuse \texttt{priority} as the ticket count.
  \item \texttt{process/process.c}, function \texttt{taskswitcher()} (around line 1315) runs on every timer interrupt, 100 times per second. It asks the current scheduler for the next process and performs the context switch with \texttt{ps\_switchto()}.
  \item \texttt{process/scheduler.c} is the default scheduler. The function \texttt{scheduler(PCB386 *lastprocess)} implements round-robin: starting from \texttt{lastprocess->next}, it returns the first process that is neither blocked nor waiting, decrementing \texttt{waiting} as it passes sleepers.
  \item \texttt{console/console.c} is the kernel shell. It already contains the \texttt{tickets} command (see below). You do \textbf{not} modify this file.
\end{itemize}

\textbf{Study \texttt{process/scheduler.c} carefully before writing any code. This is where you will place all your code. You will not change other files.}

\subsection*{Lottery scheduling}
Lottery scheduling gives each process a number of \textit{tickets}. On every scheduling decision the kernel holds a lottery:

\begin{enumerate}[itemsep=0pt,parsep=0pt]
  \item Sum the tickets of all \textit{runnable} processes. Call the sum \texttt{totaltickets}.
  \item Draw a uniform random number \texttt{winner} in $[0,\ \texttt{totaltickets})$.
  \item Walk the ready queue, accumulating the tickets of each runnable process into a counter. The first process for which \texttt{counter > winner} wins the CPU for this quantum.
\end{enumerate}

Over many quanta, a process holding three times the tickets of another receives about three times the CPU time.


\subsection*{Random number generation (provided)}
the ICS-OS kernel has no \texttt{rand()}. You must implement your own pseudo-random number generator. A \textit{xorshift32} generator (Marsaglia, 2003) is already provided for you:

\begin{minted}[frame=single,framesep=10pt]{C}
//seed/state for the pseudo-random number generator
static DWORD lottery_state = 2463534242UL;

//xorshift32 pseudo-random number generator (Marsaglia, 2003).
//The kernel has no rand(), so the scheduler brings its own.
static DWORD lottery_rand(){
	lottery_state ^= lottery_state << 13;
	lottery_state ^= lottery_state >> 17;
	lottery_state ^= lottery_state << 5;
	return lottery_state;
};

\end{minted}

\subsection*{The \texttt{tickets} command (provided)}
The kernel shell in your starter code already has a command that sets the tickets of a process:


\begin{minted}[frame=single,framesep=10pt]{bash}
tickets <pid> <ntickets>
\end{minted}

Here's the implementation in \texttt{kernel/console/console.c}:

\begin{minted}[frame=single,framesep=10pt]{C}
if (strcmp(u,"tickets") == 0){      //-- Sets the lottery tickets of a process. Args: <pid> <tickets>
	char *p1 = strtok(0," ");
	char *p2 = strtok(0," ");
	if (p1!=0 && p2!=0){
		PCB386 *proc = ps_findprocess(atoi(p1));
		if (proc != (PCB386*)-1){
			proc->priority = atoi(p2);
			printf("tickets: pid %d now holds %d ticket(s).\n", atoi(p1), atoi(p2));
		}else
			printf("tickets: process %d not found.\n", atoi(p1));
	}else
		printf("usage: tickets <pid> <ntickets>\n");
}
\end{minted}

It finds the process with \texttt{ps\_findprocess(pid)} and stores \texttt{ntickets} in its \texttt{priority} field, the same field your scheduler will read. It prints a usage message when an argument is missing and an error when the pid does not exist. \textbf{Note that we used the priority field of the PCB as the placeholder for the tickets assigned to a process}.


\subsection*{Returning the number of tickets held by a process (provided)}
This is a helper function that returns the number of tickets held by a process. For processes that were not allocated tickets (such as the startup processes and kernel threads), a default number of tickets is returned. Call this when implementing the lottery scheduler.


\begin{minted}[frame=single,framesep=10pt]{C}
//Default number of tickets given to a process that has not been
//assigned any (i.e. priority == 0)
#define DEFAULT_TICKETS 10

//returns the number of lottery tickets held by a process.
//The unused PCB386.priority field is reused as the ticket count.
static DWORD lottery_tickets(PCB386 *ptr){
	if (ptr->priority == 0) 
		return DEFAULT_TICKETS;
	return ptr->priority;
};
\end{minted}


\subsection*{The measuring tool: \texttt{lotto}}
The user application \texttt{contrib/lotto} is provided for experimental evaluation. Do \textbf{not} modify it. Since the code is quite long, it is not included in this text. It is available in the given starter code.

\begin{minted}[frame=single,framesep=10pt]{bash}
lotto.exe <ticketsA> <ticketsB> [seconds]
\end{minted}

\texttt{lotto.exe} starts two CPU-bound threads, A and B. It gives them the tickets you ask for by running the \texttt{tickets} command on them, lets them compete for \texttt{[seconds]} (default 10), and then prints how much work each one completed and what share of the total that is. If your scheduler is correct, the share of A is close to $\texttt{ticketsA} / (\texttt{ticketsA} + \texttt{ticketsB})$.

While \texttt{lotto.exe} runs, the shell is busy and does not accept typing. Wait for the prompt to come back.

The ``seconds'' are kernel seconds. Under an emulator such as QEMU, the kernel clock can run slower than your wall clock, so a 10-second run may take noticeably longer to finish.

\section{Tasks}

\subsection*{Task 1: Code Tracing (15 points)}
No coding is needed here. Read the source files described in the Background and answer the questions.

\qbox{
\textbf{QUESTIONS}(Provide screen shots to support your answer):
\begin{enumerate}[itemsep=0pt,parsep=0pt]
  \item{Who calls \texttt{scheduler()}, and how often? (Hint: follow \texttt{taskswitcher()} and the timer interrupt.)}
  \item{Why does the round-robin loop decrement \texttt{ptr->waiting} while scanning? What would happen to sleeping processes if it did not?}
  \item{The comment says the scheduler ``assumes that at least one process is not blocked or waiting.'' Why is this assumption safe in ICS-OS? (Hint: what do the kernel process and its threads do?)}
\end{enumerate}
}

\subsection*{Task 2: Baseline, Measure Round-Robin (5 points)}
Before you change anything, build the stock kernel (one with default scheduler) and measure it. Run the commands below (the first group inside the Docker build environment).

\begin{minted}[frame=single,framesep=10pt]{bash}
$docker compose --remove-orphans run ics-os-build
(in container)# cd /home/ics-os && make clean && make
(in container)# cd contrib/lotto && make && make install && cd ../..
(in container)# exit
$sudo make floppy
$qemu-system-i386 -fda ics-os-floppy.img
\end{minted}

At the ICS-OS prompt run the two commands below.

\begin{minted}[frame=single,framesep=10pt]{bash}
[1]% lotto.exe 10 10 20
[2]% lotto.exe 30 10 20
\end{minted}

\texttt{lotto.exe} will report that thread A and thread B hold the tickets you asked for, but the stock scheduler ignores them. Capture a screen shot and record the shares printed. Round-robin ignores tickets, so \textbf{both} runs should split the work close to 50\% / 50\%, even though the second run gives A three times the tickets of B.

\subsection*{Task 3: Implement Lottery Scheduling (35 points)}
Edit \texttt{process/scheduler.c}. This is the only file you modify. Replace the round-robin logic in \texttt{scheduler()} with lottery scheduling:

\begin{enumerate}[itemsep=0pt,parsep=0pt]
  \item Add a \texttt{lottery\_rand()} xorshift32 PRNG (a \texttt{static DWORD} state variable with a non-zero seed).
  \item Add a helper \texttt{lottery\_tickets(PCB386 *ptr)} that returns the process's ticket count. Use the unused \texttt{priority} field as the ticket count. If it is \texttt{0} (the default for every process), return \texttt{DEFAULT\_TICKETS} (use 10).
  \item \textbf{Pass 1:} traverse the circular list exactly once, starting at \texttt{lastprocess->next}. Decrement \texttt{waiting} on sleepers, and sum the tickets of runnable processes (not \texttt{PS\_ATTB\_BLOCKED}, not \texttt{waiting}).
  \item If \texttt{totaltickets == 0}, fall back to the original round-robin loop, so that the kernel still behaves if nothing is runnable on this tick.
  \item \textbf{Pass 2:} draw \texttt{winner = lottery\_rand() \% totaltickets}, walk the list again accumulating the tickets of runnable processes, and return the winner.
  \item Update the strings in \texttt{ps\_scheduler\_install()} (\texttt{hdr.name}, \texttt{hdr.description}, and the startup \texttt{printf}) so that they say lottery instead of round-robin.
\end{enumerate}

\noindent\textbf{Constraints.} Do \textbf{not} modify \texttt{PCB386}, \texttt{taskswitcher()} or \texttt{console.c}. Your whole change lives in \texttt{scheduler.c}. Interrupts are already disabled when \texttt{scheduler()} runs, so no extra locking is needed. Your function must terminate, and it must never return a blocked or waiting process.

\subsection*{Task 4: Build and Experiment (20 points)}
Rebuild the kernel and the floppy image as in Task 2 and boot it. Only \texttt{scheduler.c} has changed, so you do not need to rebuild \texttt{lotto}.

\begin{enumerate}[itemsep=0pt,parsep=0pt]
  \item The system must boot to the shell and \texttt{ps} must work. If your scheduler is wrong, the system usually hangs during boot or at the first command.
  \item Run the three commands below. Run each at least \textbf{three times}.
\begin{minted}[frame=single,framesep=10pt]{bash}
[1]% lotto.exe 10 10 20
[2]% lotto.exe 30 10 20
[3]% lotto.exe 10 40 20
\end{minted}
  \item Record the share that each run reports and compare it with the share the program lists as expected.
  \item Put your results in a table with these columns: tickets, expected share, measured share for each trial, and the average. Add the Task 2 baseline results for comparison.
\end{enumerate}

Here is a sample of what a short run (\texttt{lotto.exe 30 10 8}) prints. Your numbers will differ.

\begin{minted}[frame=single,framesep=10pt]{text}
lotto: thread A = pid 21, thread B = pid 22
tickets: pid 21 now holds 30 ticket(s).
tickets: pid 22 now holds 10 ticket(s).
lotto: running for 8 seconds...
tickets   A = 30   B = 10
work      A = 460331   B = 172394   (units of 1000 loops)
share     A = 72%   B = 27%
expected  A = 75%   B = 25%
\end{minted}

The two shares can add up to 99\% because each is rounded down.

\section{Report Questions (10 points)}

\qbox{
\textbf{QUESTIONS}:
\begin{enumerate}[itemsep=0pt,parsep=0pt]
  \item{Your measured shares will not be exactly the expected ones. Give two reasons. (Think: how many lotteries are held in a 20-second run at 100 ticks per second, and which other processes hold tickets?)}
  \item{The baseline in Task 2 ignored tickets, and the lottery kernel did not. Which lines of your \texttt{scheduler()} are responsible for the difference?}
  \item{Lottery scheduling is probabilistic. Describe \textit{stride scheduling}, its deterministic counterpart, and one advantage it has.}
  \item{What happens if a malicious user program could set its own tickets to a huge value? How do real systems guard proportional-share mechanisms?}
  \item{Why is it acceptable for the scheduler to use a weak, non-cryptographic PRNG here?}
\end{enumerate}
}

\section{Reflection}
Write some realizations and questions that crossed your mind while doing this lab.

\section*{Reference}
C. A. Waldspurger and W. E. Weihl. \textit{Lottery Scheduling: Flexible Proportional-Share Resource Management.} OSDI 1994.

\end{document}
